% ======================================================================
%  05_branes.tex —— 第 5 讲 膜、线算子、特征膜
%  依据：Kapustin--Witten 2006 \S6--\S8（印刷 pp.78--123）
% ======================================================================

\section{膜、线算子、特征膜}
\label{sec:branes}

\begin{goals}
\begin{itemize}
\item 说清 2d sigma 模型的边界条件为什么就是膜，A-膜与 B-膜各保留哪个超荷、各要什么样的几何（Lagrange/余各向同性 对 复子流形）。
\item 掌握 Wilson 算子与 't~Hooft 算子的定义差别（order 对 disorder），以及它们各自出现在哪个 $\Psi$ 值。
\item 理解特征膜的定义 $XB\cong B\otimes V$：为什么「特征值」必须换成向量空间，为什么零膜是 Wilson 的特征膜、Hitchin 纤维是 't~Hooft 的特征膜，以及全体特征膜如何合成 $C$ 上一条局部系统。
\end{itemize}
\end{goals}

前几讲已经把舞台搭好：四维 $\mathcal N=4$ 超杨--米尔斯理论紧化在曲线 $C$ 上，低能有效理论是一个以 $\MH$ 为靶空间的二维 sigma 模型，$\MH$ 超凯勒、带三个复结构 $I,J,K$（\src{B p.60}）。本讲放三样东西：膜、线算子、特征膜（内容据原文 \S6--\S8，\cite{kw06}）。先说膜分几类；再说 Wilson 与 't~Hooft 线算子如何降成二维算子；最后说算子作用在膜上何时「本征」，以及全体特征膜怎样拼出 $C$ 上的局部系统——那是几何 Langlands 的伽罗瓦一侧在物理里的第一次显形。

\subsection{边界条件就是膜：A 型与 B 型}

二维拓扑场论最容易被忽略的一点是：有边界时，理论的一半数据在边界上。世界面 $\Sigma$ 带边界时，边界条件要指定三样东西：场 $\Phi:\Sigma\to X$ 在 $\partial\Sigma$ 上落在哪个子流形 $N\subset X$ 里；边界上带什么 Chan--Paton 丛与联络；法向导数满足什么一阶条件。超对称把这些数据卡得很死：要让边界保留某一个超荷 $Q$，$(N,\text{丛})$ 就必须满足一条几何方程。选不同的 $Q$，就得到不同类型的膜（\src{B p.77}--\src{B p.78}）。

先看 A 型。A 型超荷的边界项正比于 $\int_{\partial\Sigma}\omega(\delta\Phi,\partial_t\Phi)$，要它为零，必须 $\omega|_N=0$，即 $N$ 是 Lagrange 子流形；同时边界规范场必须平坦，$F=0$。把 Lagrange 放宽成\term{coisotropic}{余各向同性}子流形（$TN^{\omega}\subset TN$）并配上横截的复结构，也能给出 A-膜；这是后文第 7 讲真正要用的情形。A 型超荷的上同调要数世界面瞬子（全纯圆盘），计算它用的是 Floer 型或量子上同调的技术。

再看 B 型。B 型超荷里含 $\bar\partial$，保它要求 $N$ 是复子流形、Chan--Paton 丛是全纯的。B 型超荷的上同调就是 Dolbeault 上同调，世界面瞬子不贡献修正。所以一句话：A 型要 Lagrange（或余各向同性）加平坦丛，保的是「辛」的那个超荷；B 型要复子流形加全纯丛，保的是「复」的那个超荷。

$\MH$ 是超凯勒的，三个复结构各有一对 A 型与 B 型超荷，因此一张膜可以「对 $I$ 是 B 型、对 $J$ 是 A 型」。真正有意思的是同时保留二维理论一半超对称的膜（这也等于保留四维理论四分之一超对称）：这样的膜只有四种基本类型，记作
\[
(B,B,B),\qquad (B,A,A),\qquad (A,B,A),\qquad (A,A,B),
\]
括号里三个字母依次表示该膜在 $I,J,K$ 中的类型（\src{B p.78}）。它们与两个常用例子一起列在表~\ref{tab:brane-types} 里。

\begin{table}[htbp]
\centering
\small
\caption{四类保留二维一半超对称的膜。括号中字母依次对应复结构 $I,J,K$；末列是典型例子（据 \src{B p.78} 与 \src{B p.72} Table~2）。}
\label{tab:brane-types}
\begin{tabular}{@{}lcccc@{}}
\toprule
类型 & $I$ & $J$ & $K$ & 典型例子 \\
\midrule
$(B,B,B)$ & B & B & B & 零膜；空间充满膜（全纯丛） \\
$(B,A,A)$ & B & A & A & Hitchin 纤维；固定 $E$ 让 $\varphi$ 变 \\
$(A,B,A)$ & A & B & A & $J$ 中的复 Lagrange 子流形 \\
$(A,A,B)$ & A & A & B & $K$ 中的复 Lagrange 子流形 \\
\bottomrule
\end{tabular}
\end{table}

两个例子要记牢。第一个是\term{zerobrane}{零膜} $B_x$：支撑在点 $x\in\MH$ 上，Chan--Paton 丛是支撑在 $x$ 的摩天楼层；点在任何复结构下都是复子流形，丛在点上全纯，故零膜是 $(B,B,B)$ 型。第二个是 Hitchin 纤维 $F$：它在复结构 $I$ 中是复子流形，并且对全纯辛形式 $\Omega_I$ 是 Lagrange 的（因而对 $\omega_J,\omega_K$ 也是），配上平凡丛，故是 $(B,A,A)$ 型（\src{B p.78}--\src{B p.79}）。

另外两条。其一，A-膜的支撑维数至少是靶空间的中维：Lagrange 最常见，余各向同性的支撑可以更高——第 7 讲构造 D-模时用的正是这个极端情形。其二，S-对偶把 $(B,B,B)$ 换成 $(B,A,A)$，把 $(A,B,A)$ 与 $(A,A,B)$ 各自映到自身（\src{B p.72}）；这个交换本讲末尾会再次出现。

\subsection{线算子：Wilson 与 't~Hooft}

四维理论里，算子可以支撑在任意子流形上。支撑是一维时叫线算子。理论是拓扑的，把支撑连续形变不改变算子，于是线算子只依赖支撑的同伦类：支撑是闭圈时叫圈算子，支撑是一条跑到无穷远的线时叫线算子。路径积分里插入线算子，相当于在它周围挖掉一个小管、在管的边界上指定某种数据；指定方式不同，算子就不同。两种最重要的方式是 Wilson 与 't~Hooft，分别出现在对称性增强的两个特殊值 $\Psi=\infty$ 与 $\Psi=0$（\src{B p.78}--\src{B p.79}）。

Wilson 算子在 $\Psi=\infty$（即 $t=\pm\ii$）。取规范群的表示 $R$，沿圈 $S$ 取 holonomy 的迹：
\begin{equation}
W_R(S)=\Tr_R\,P\exp\Big(-\oint_S\mathcal A\Big),\qquad \mathcal A=A+\ii\varphi,
\label{eq:wilson}
\end{equation}
其中 $A$ 是规范场，$\varphi$ 是复标量。为什么只用这个组合：在 $t=\pm\ii$ 处，拓扑扭化的变换把 $A$ 与 $\varphi$ 一起旋转，唯一保持不变的复化联络就是 $A+\ii\varphi$。为什么它是拓扑的：在 $t=\pm\ii$，BPS 方程蕴含 $\mathcal A$ 平坦，于是 holonomy 只依赖 $S$ 的同伦类，把 $S$ 挪一挪不改变 $W_R(S)$。开线也不浪费：从 $p$ 到 $q$ 的开线给出一个平行移动矩阵，可以在两端之间搬运表示（\src{B p.79}--\src{B p.80}）。

't~Hooft 算子在 $\Psi=0$（即 $t=1$、$\theta=0$）。它不是某个联络的迹，而是一个 disorder 算子：在支撑 $S$ 附近指定经典场的奇点，然后对「带这个奇点的场」做路径积分——算子本身就是这个受限制的路径积分。$G=\mathrm{U}(1)$ 时，磁荷 $m$ 的奇点就是 Dirac 单极，绕 $p$ 的小球面上 $\oint_{S^2}c_1(L)=m$（\src{B p.83}）。一般 $G$：奇点由同态 $\rho:\mathrm{U}(1)\to G$ 的共轭类指定；GNO 对偶说，这样的共轭类等价于 $\Ldual$ 的一个最高权，也等价于 $\Ldual$ 的一个不可约表示 ${}^{L}R$，于是把 't~Hooft 算子记作 $T({}^{L}R,S)$。这里还有一层离散数据：$\xi({}^{L}R)\in H^2(S^2,\pi_1(G))$ 度量算子的扭曲程度，它非平凡当且仅当 $\Ldual$ 的中心在 ${}^{L}R$ 上作用非平凡（\src{B p.84}）。

两个算子的对比值得背下来：Wilson 是 order 算子，它读出联络；'t~Hooft 是 disorder 算子，它指定奇点，算子的值是奇点存在下的路径积分。S-对偶 $\Psi\mapsto-1/(n_g\Psi)$ 把 $\Psi=\infty$ 与 $\Psi=0$ 互换，于是把 Wilson 与 't~Hooft 互换，并把表示 $R$ 换成 ${}^{L}R$（\src{B p.98}）。这个互换是本讲后半段与下一讲的引擎。

\subsection{线算子作用在膜上：特征膜}

现在让线算子作用在膜上。做法是「折叠」：线算子支撑在 $\Sigma$ 上的一条线（再乘上 $C$ 上的点 $p$）；把 $\Sigma$ 沿这条线剪开再翻折，得到一个界面——它是两套二维理论的公共边界条件，当然也是膜，准确说是一个 $(Y,Z)$-膜。把线算子 $X$ 变成界面，这一步就是把算子「打开」成一个边界条件。膜 $B$ 与界面复合，结果仍是膜，记作 $XB$：线算子把膜映到膜（\src{B p.90}）。

用态空间说更准确：设 $H_{B_1,B_2}$ 是两端带边界条件 $B_1,B_2$ 的区间上的态空间（边界算子的空间），则 $XB$ 的作用由 $H_{\tilde B,XB}$ 刻画；$H_{B,B}$ 是结合代数（边界算子的 OPE），$H_{B_1,B_2}$ 是双模（\src{B p.91}--\src{B p.92}）。

接下来定义本讲的主角。普通的特征方程 $O\psi=\lambda\psi$ 在这里没有意义：算子作用的对象是膜（一个范畴的对象），不是向量，没有「本征值」可谈。要类比，就得把 $\lambda$ 换成同样能「张量上去」的东西——一个向量空间。先定义膜与向量空间的张量积 $B\otimes V$：几何上就是把 Chan--Paton 丛 $W$ 换成 $W\otimes V$（$V=\CC^k$ 时就是取 $k$ 份膜）；抽象地由
\begin{equation}
H_{\tilde B,B\otimes V}=H_{\tilde B,B}\otimes V
\label{eq:tensor}
\end{equation}
刻画。若
\begin{equation}
XB\cong B\otimes V,
\label{eq:eigenbrane}
\end{equation}
就说 $B$ 是 $X$ 的\term{eigenbrane}{特征膜}，$V$ 是它的「特征值」（\src{B p.95}）。

这个定义有两点要紧。第一，特征值不是数而是向量空间；一旦有了向量空间，「特征值随参数变化」才有意义。第二，式~\eqref{eq:eigenbrane} 正是 Beilinson--Drinfeld 在几何 Langlands 里定义的 eigensheaf 的物理版；这不是巧合，而是第 6 讲要还的债。

第一个例子（电）：Wilson 算子的特征膜是零膜。Wilson 算子 $W_q(R)$ 作用在膜上的方式，是把它的 Chan--Paton 丛 $U$ 换成 $U\otimes E_q(R)$，其中 $E_q(R)$ 是从 $C$ 上的通用丛诱导、在表示 $R$ 下取纤维的丛。零膜 $B_x$ 的 Chan--Paton 丛是支撑在 $x$ 的摩天楼层 $U_x$；任何丛限制到一个点都平凡，于是
\begin{equation}
U_x\otimes E_q(R)=U_x\otimes E_q(R)|_x,
\label{eq:zerobrane}
\end{equation}
右边只是 $x$ 处的一个固定向量空间。故 $W_q(R)B_x\cong B_x\otimes E_q(R)|_x$：零膜是电特征膜，特征值是「表示空间在 $x$ 处的纤维」（\src{B p.121}）。

第二个例子（磁）：'t~Hooft 算子的特征膜是 Hitchin 纤维，准确说是「卷在纤维上、带秩 1 平坦 Chan--Paton 丛」的膜，原文叫 type~F 膜。它和零膜的关系正是 S-对偶：$\MH(\Ldual,C)$ 里的零膜，经 S-对偶变成 $\MH(G,C)$ 里 Hitchin 纤维上的 type~F 膜。为什么非要从点换成纤维？因为电与磁「局域化」的方向相反：零膜把膜局域在一个点，读出电表示；磁荷的信息沿 Hitchin 纤维的环面方向铺开，type~F 膜正好卷在整条纤维上，它的平坦 Chan--Paton 丛给出对偶环面上的一个点，这个点就是 't~Hooft 算子的「特征值」。原文还用 OPE 印证这件事：Wilson 算子的 OPE 生成 $G$ 的表示环（原文 (6.27)），而 't~Hooft 算子的 OPE 应该生成 $\Ldual$ 的表示环（原文 (6.29)；原文在这里写的是 “we expect”）（\src{B p.98}--\src{B p.99}）。

还有一条结构性事实：从四维来的线算子自动两两交换。物理理由很直接：它们支撑在 $C$ 上的不同点，四维里根本不相交，可以随意前后移动；$C$ 是二维的，两个不同点在空间中的位置足够挪开而不产生单值性。既然算子交换，就可以同时找它们的特征膜——这是下一讲把整个算子族放在一起讨论的前提（\src{B p.97}）。

\subsection{全体特征膜 = 曲线上的局部系统}

最后把所有特征膜的「特征值」沿 $C$ 拼起来。线算子可以放在任意 $p\in C$，记 $X_p$。若 $B$ 对某个 $p$ 是特征膜，则对所有 $p$ 都是：
\begin{equation}
X_pB\cong B\otimes V_p,\qquad p\in C .
\label{eq:family}
\end{equation}
单看一个 $p$，$V_p$ 是一个向量空间；但 $p$ 可以连续移动。移动 $p$ 相当于改变 $C$ 的度规，作用量只变一个 $Q$-恰当项，故可在族 $V_p$ 上装一个平坦联络，使 $V_p$ 局部与 $p$ 无关；整体上可以有非平凡单值性。用几何语言：$p\mapsto V_p$ 是 $C$ 上一个平坦向量丛（\src{B p.96}）。

不同表示的数据还要合起来。OPE 给出
\begin{equation}
W(R)W(\tilde R)=\sum_\alpha n_\alpha\,W(R_\alpha),\qquad R\otimes\tilde R=\bigoplus_\alpha n_\alpha R_\alpha,
\label{eq:wilson-ope}
\end{equation}
特征值必须满足同样的分解，于是 $V_p(R)\otimes V_p(\tilde R)\cong V_p(R\otimes\tilde R)$；用范畴的话说，特征膜给出一个张量函子 $\mathrm{Rep}(G)\to\mathrm{Vect}$。Tannaka 对偶说，这样的张量函子等价于一个局部系统。电侧（零膜）得到的是取值在 $G_{\CC}$ 里的局部系统；磁侧（type~F 膜）得到的是取值在 $\LdualC$ 里的局部系统
\begin{equation}
\rho:\pi_1(C)\longrightarrow \LdualC .
\label{eq:local-system}
\end{equation}
这就是「全体特征膜合起来等于 $C$ 上一条局部系统」的准确含义（\src{B p.96}--\src{B p.98}）。

为什么说这已经像几何 Langlands？几何 Langlands 的标准故事正好从这里开始。给定 $\Ldual$ 与同态 $\vartheta:\pi_1(C)\to\LdualC$（并要求它半稳定），得到 $\MH(\Ldual,C)$ 中一点 $x(\vartheta)$；$\vartheta$ 不可约时该点光滑，其零膜是电特征膜；S-对偶把它变成 $\MH(G,C)$ 中 Hitchin 纤维上的 type~F 磁特征膜（\src{B p.122}）。几何 Langlands 断言，这个 $\vartheta$ 应该对应 $\MC$ 上一个 Hecke 特征层——既是 Hecke 算子的特征对象，又是 D-模。物理侧现在有了「磁特征膜」这一候选对象，但还没有证明它就是 Hecke 特征层；这一步正是第 6 讲的任务。原文在这里的措辞很克制：它只说我们「至少可以看到一线微光」（at least a glimmering）（\src{B p.122}--\src{B p.123}）。

所以本讲的收获是：线算子把膜映到膜，只乘一个向量空间的膜是特征膜，两类特征膜的特征值沿 $C$ 拼成局部系统。语言一旦立住，剩下的问题就只有一个：'t~Hooft 算子到底是不是数学上的 Hecke 算子？

\begin{keybox}
线算子把膜映到膜；特征膜是满足 $XB\cong B\otimes V$ 的膜，「特征值」是向量空间而非数；电侧特征膜是零膜，磁侧是 Hitchin 纤维上的 type~F 膜，其特征值沿 $C$ 拼成局部系统——几何 Langlands 的伽罗瓦一侧在物理里第一次显形。
\end{keybox}

\begin{follow}
\begin{itemize}
\item B \S6, pp.78--102：线算子的定义（(6.4)--(6.5)）、't~Hooft 的奇点与 GNO 对偶（(6.9) 附近）、折叠与特征膜、OPE。第一遍抓住 (6.22) 与 (6.27) 即可。
\item B \S7, pp.103--115：中心、离散 flux 与 twisted 丛，可跳读，只记住「电/磁 flux 成对、S 交换它们」。
\item B \S8, pp.116--123：零膜 $=$ 电特征膜（(8.10)）、type~F 膜 $=$ 磁特征膜、原文对几何 Langlands 的「一线微光」；这是第 6 讲的直接入口。
\end{itemize}
\end{follow}

\begin{checkbox}
\begin{enumerate}
\item 为什么 A-膜要求 $\omega|_N=0$（或余各向同性）而 B-膜要求 $N$ 是复子流形？它们各自保留哪个超荷，对应的上同调分别是什么？
\item Wilson 算子为什么只依赖圈的「同伦类」？'t~Hooft 算子为什么不是某个联络的 holonomy，而是一个 disorder 算子？
\item 特征膜的定义里为什么必须用向量空间代替数？写出 $XB\cong B\otimes V$ 的态空间刻画，并说明 $B\otimes V$ 在几何上到底是什么操作。
\item 零膜为什么是 Wilson 算子的特征膜？把「任何丛限制到点都平凡」这一步补全。
\item Hitchin 纤维为什么是 't~Hooft 算子的特征膜？它和零膜通过什么操作联系起来？为什么磁侧非要从点换成纤维？
\item 电侧、磁侧两组特征膜分别给出什么群里的局部系统？几何 Langlands 用的是哪一侧，为什么？
\end{enumerate}
\end{checkbox}
